\begin{table}[htbp]
\centering
\small
\caption{Candidate-park evidence states in the 19 November 2026 OpenStreetMap
snapshot. No candidate has verified public access or a verified physical
entrance. All 1,015 candidates remain unresolved.}
\label{tab:park}
\begin{threeparttable}
\begin{tabular}{@{}lrr@{}}
\toprule
Evidence state & Candidates & Share \\
\midrule
Unknown access & 980 & 96.55\% \\
Conditional, unverified & 28 & 2.76\% \\
Positive OSM tag, unverified & 7 & 0.69\% \\
Verified public access & 0 & 0.00\% \\
Verified physical entrance & 0 & 0.00\% \\
Unresolved review status & 1,015 & 100.00\% \\
\bottomrule
\end{tabular}
\begin{tablenotes}[flushleft]\footnotesize
\item Evidence states are hierarchical and are not simple access-tag counts. A positive OSM tag is an unverified assertion; an unknown tag is not evidence of closure. The positive-tag-only scenario retains 10 destinations and is a coverage stress test, not a better estimate of public access.
\end{tablenotes}
\end{threeparttable}
\end{table}
